% 03_restart.tex - split a run in two and continue it from solution.rvc
\section{Continue a run from \texttt{solution.rvc}}
At the end of every run, Raven writes the state of the model to \cmd{solution.rvc}: the soil and surface stores, the
channels and the reservoirs, and in a coupled run the head of every MODFLOW cell (and, with a soil zone, its columns). Used
as the \cmd{.rvc} file (the initial conditions) of the next run, it continues the run where it stopped.

\step{Copy the start model and make a continuous run}
The reference is the three years 2010--2012 in one run:
\begin{shell}
$ cp -r $WS/model/start $WS/work
$ cd $WS/work
\end{shell}
\term{t5_run_full.txt}
\see \cmd{Exiting Gracefully: Successful Simulation}.

\step{Run part 1: the first two years}
Change the duration in \cmd{Creek.rvi} to two years and run into \cmd{out\_part1}:
\diff{t5_part1.diff}
\term{t5_run_part1.txt}
\see \cmd{Exiting Gracefully: Successful Simulation}.

\step{Look at the saved state}
\term{t5_rvc_blocks.txt}
\why after Raven's own blocks (HRU states, basin states) comes \cmd{:GWHeads \V{t5:gwheads}}: the head of every cell,
layer by layer and row by row, with 17 significant digits:
\term{t5_rvc_head.txt}

\step{Make a folder for part 2}
Make a folder for the second part with the model files, and use \cmd{solution.rvc} of part~1 as its \cmd{.rvc} file:
\begin{shell}
$ mkdir ../work_part2
$ cp Creek.* CreekMet.rvt CreekGW.rvt *.geojson *.tif ../work_part2/
$ cp out_part1/solution.rvc ../work_part2/Creek.rvc
$ cd ../work_part2
\end{shell}
\see no output.

\step{Run part 2: the third year, from the hotstart}
Change \cmd{Creek.rvi} in \cmd{work\_part2} to start on 1 January 2012 and to run the rest of the three years:
\diff{t5_part2.diff}
\term{t5_run_part2.txt}
\term{t5_part2_advisory.txt}
\see \cmd{Exiting Gracefully: Successful Simulation}, and a note that the groundwater heads come from the hotstart.
\why the line \cmd{:GWInitialization STEADY\_STATE} of the \cmd{.rvg} is skipped: the heads come from the hotstart file.

\step{Compare with the continuous run}
\term{t5_diff.txt}
\begin{tocheck}
The largest difference is \V{t5:maxdiff}: the restarted run reproduces the continuous run bit for bit.
\end{tocheck}
\begin{trouble}
\begin{itemize}
\item \emph{``running more time steps than the MODFLOW 6 simulation holds''}: \cmd{:Duration} or \cmd{:TimeStep}
  changed after Raven wrote the MODFLOW files. Run Raven again into a new output folder.
\item A damaged \cmd{:GWHeads} header or a non-number in a groundwater block stops the run with a message that names it.
  Copy \cmd{solution.rvc} again.
\end{itemize}
\end{trouble}
\why restarts have three common uses. Spin a model up once (run it until its stores stop trending up or down) and start every later run from that state. Split long runs
into parts. Carry the water in the unsaturated columns of the soil-zone coupling into the next run: it is saved in
\cmd{solution.rvc} as well.
